\section{Introduction}\label{section: introduction}

The \emph{star height} of a regular expression is the maximal nesting depth of Kleene stars in it~\cite{egganTransitionGraphsStarheight1963}.
When complement is admitted as an operator, one speaks of \emph{generalized} star height.
Languages of generalized star height $0$ are exactly the star-free languages~\cite{schutzenbergerFiniteMonoidsHaving1965}.
It is a long-standing open problem whether some regular language has generalized star height at least two~\cite{pinResultsGeneralizedStarheight1992}.
Recently, OpenAI announced a proof that every regular language has generalized star height at most three~\cite{openaiGeneralizedStarHeight2026}.

The same question can be asked for \emph{terms} of Kleene algebra with complement, that is, for expressions over variables.
Here the variables are interpreted as binary relations, and two terms are compared by their values under all valuations in a class of models.
This viewpoint goes back to Tarski's calculus of relations~\cite{tarskiCalculusRelations1941} and its extension with transitive closure~\cite{ngRelationAlgebrasTransitive1984}.
We work with $\RSUB$~\cite{nakamuraFiniteRelationalSemantics2025}.
In $\RSUB$, the universe is a finite linear order, each variable denotes a set of forward pairs, and complement is taken relative to the set of all forward pairs.
$\RSUB$ captures the equational theory of Kleene algebra with complement with respect to languages~\cite{nakamuraFiniteRelationalSemantics2025}.

In this paper we ask \emph{how simple the structures can be} if generalized star height one and two are still to be separated.
We measure structural simplicity by pathwidth.
To a valuation we associate its Gaifman graph, in which we also include the successor relation of the underlying linear order.
The successor relation is star-free definable from the order, so this is a natural choice.
Let $\Cls{k}$ be the class of $\RSUB$ valuations whose graph has pathwidth at most $k$.
\begin{itemize}
  \item On $\Cls{1}$, every input pair is a loop or a successor pair, since any other pair closes a cycle with the successor path.
        So a valuation in $\Cls{1}$ is essentially a word whose letters record which variables hold on each successor pair.
        The separation of heights one and two on $\Cls{1}$ is therefore closely related to the classical open problem (\Cref{remark: pathwidth one words}).
  \item Our main result shows that one more unit of pathwidth suffices for a separation.
\end{itemize}

\begin{theorem}[Main theorem]\label{theorem: main}
  There is a term $\Psit$ over three variables with star height two such that no term of star height at most one is equivalent to $\Psit$ on $\Cls{2}$.
  In other words, the generalized star height of $\Psit$ relative to $\Cls{2}$ is exactly two.
\end{theorem}

The term $\Psit$ is defined in \eqref{equation: Psi}.
It is fixed in advance, and only the counter-models depend on the competing term.
The counter-models have a very simple shape (\Cref{figure: module,figure: cells}).
Their graphs consist of a single path, together with a few vertices that are joined to some later vertices of their own segment of the path.
So each segment forms a fan, and the whole graph is a chain of fans.
Such graphs are outerplanar and have pathwidth exactly two.
The bound two is tight for this witness.
On $\Cls{1}$, the term $\Psit$ is equivalent to a term of star height one (\Cref{proposition: pathwidth one}).

\subparagraph*{Idea of the proof.}
The term $\Psit$ has the form $\Ht^*$.
Here $\Ht$ is a star-free combination of a relation $\Kt$ and its complement $\Nt$ inside a fixed relation $\vb$, and $\Kt$ contains the star $(\vr \compo \vr)^*$.
The counter-models are sequences of \emph{cells}.
Each cell contains a \emph{module}, that is, an interval $[u, v]$ with $(u, v) \in \vb$.
The inner star decides whether $(u, v)$ belongs to $\Kt$ by reading the parity of the length of an $\vr$-chain.
The outer star composes the cells, and the cells act as the identity or as a swap on two ports.
So $\Psit$ computes the parity of the number of cells whose module lies in~$\Kt$.

For a competing term $E$ of star height one, let $q$ bound the composition depth of the bodies of its stars.
We construct two modules $\M^{+}_q$ and $\M^{-}_q$ on the same set of vertices (\Cref{proposition: local pair}).
They differ in a single input pair, and $(u, v) \in \Kt$ holds in $\M^{+}_q$ but not in $\M^{-}_q$.
Still, no star of depth at most $q$ can tell them apart.
To construct them, we encode prefixes of a module by first-order types and suffixes by a finite automaton.
The first-order types form an aperiodic monoid, while the automaton may count modulo a number.
An idempotent in the minimal ideal of the product monoid lets us rewind both components at the same time (\Cref{section: local}).
A replacement lemma then shows that every star in $E$ is blind to the modules strictly between its endpoints (\Cref{section: replacement}).
Hence $E$ translates into first-order logic with modular predicates over the sequence of cell bits.
That logic cannot express parity, so $E$ differs from $\Psit$ (\Cref{section: lower bound}).

\subparagraph*{Related work.}
For regular languages, the generalized star-height problem is discussed in~\cite{pinResultsGeneralizedStarheight1992}.
Inside generalized star height zero, the dot-depth hierarchy is infinite~\cite{brzozowskiDotdepthHierarchyStarfree1978,thomasConcatenationGameDotdepth1987}.
% TODO: cite the companion result on the star-height hierarchy of terms w.r.t. languages / RSUB, if appropriate.
Terms of Kleene algebra with complement can be read as formulas of transitive closure logic with three variables.
For transitive closure logic on finite structures, Gr\"adel and McColm proved hierarchy theorems for the nesting of transitive closure and negation~\cite{gradelHierarchiesTransitiveClosure1996}.
On ordered finite structures, every formula of transitive closure logic is equivalent to one with a single transitive closure, at the price of higher arity~\cite{immermanLanguagesThatCapture1987}.
So the bounded arity of terms is essential here.
Expressive-power results for relational calculi via Ehrenfeucht--Fra\"iss\'e games include~\cite{nakamuraExpressivePowerSuccinctness2022}.
Our use of an idempotent from the minimal ideal of a finite monoid is standard in the algebraic theory of regular languages~\cite{pinVarietiesFormalLanguages1986}.

\subparagraph*{Organization.}
\Cref{section: preliminaries} fixes notation.
\Cref{section: overview} defines $\Psit$ and outlines the proof of \Cref{theorem: main}.
\Cref{section: modules,section: local,section: cells,section: replacement,section: lower bound} contain the proof.
\Cref{section: pathwidth one} treats pathwidth one, and \Cref{section: conclusion} concludes.
